%% RESUMO NA LINGUA VERNACULA (elemento obrigatorio -- NBR 6028)
%% Versao dummy do exercicio ia-6.1 (o resumo final tera de 150 a 500 palavras).

\begin{resumo}
Esta dissertação investiga recomendações de investimento em renda fixa geradas
com apoio de modelos de linguagem para o investidor pessoa física brasileiro,
no recorte de Tesouro Direto, CDB, LCI e LCA. O objetivo é medir duas
propriedades de cada recomendação: a adequação do produto ao perfil declarado
do investidor e a fidedignidade da justificativa em linguagem natural, isto é,
a ausência de afirmações que contrariem a ficha do produto. Perfis sintéticos
são submetidos a um protótipo de recomendador, e cada justificativa é comparada
com a ficha do produto segundo uma tipologia de alucinação intrínseca e
extrínseca. Espera-se identificar em que condições a justificativa gerada pode
ser oferecida ao investidor sem revisão humana caso a caso.
\end{resumo}
